\section{Open Questions}

The following are genuinely open questions raised by (but not answered by)
this replication. Each is accompanied by a concrete next-step probe.

\begin{enumerate}
\item \textbf{Do the QWC/FC grouping benefits (QWC $\sim$3$\times$,
  FC $\sim$50$\times$ on LiH) survive at realistic VQE trial states rather
  than the exact ground state?}
  \emph{Basis.} Our LiH numbers are ground-state variances via dense
  diagonalization; per-fragment variances differ near approximate VQE minima,
  and the review does not report ansatz-conditioned variance surfaces for
  our molecules.
  \emph{Next steps.} Re-run a short UCCSD-VQE loop (Qiskit Nature +
  statevector) on LiH; record per-fragment variances at each iterate and
  the converged $\theta^*$; recompute $M_{\mathrm{opt}}$ at $\theta^*$ for
  QWC and FC; report the ratio $M_{\mathrm{opt}}(\theta^*)/M_{\mathrm{opt}}(\text{gs})$.
  Repeat with a hardware-efficient ansatz.

\item \textbf{Grouping vs.\ modern classical-shadow variants on the same
  Hamiltonian.}
  \emph{Basis.} The review discusses both grouping and shadows but does not
  provide a like-for-like head-to-head at a fixed molecule and shot budget.
  Our baseline is naive single-Pauli-per-shot; a well-tuned locally-biased
  shadow protocol could outperform FC grouping.
  \emph{Next steps.} Implement locally-biased classical shadows
  (Huang--Kueng--Preskill 2020) plus Hadfield-style derandomization (2022)
  for LiH; hold total shots fixed; measure realized MSE over 200 repeats at
  1.5k / 15k / 150k shots; compare directly to the FC-grouping numbers in
  \texttt{grouping\_summary\_v3.json}.

\item \textbf{Does FC$>$QWC hold for strongly-correlated systems (stretched
  N$_2$, Cr$_2$, small transition-metal complex)?}
  \emph{Basis.} H$_2$ and equilibrium LiH are weakly / mildly correlated;
  strongly-correlated states have heavier-tailed Pauli-coefficient
  distributions that can change the greedy-coloring landscape. The review
  asserts FC$>$QWC generically but does not stress-test strong correlation.
  \emph{Next steps.} Build stretched-N$_2$ (R=2.5\,\AA{}, STO-3G or
  cc-pVDZ frozen-core) and Cr$_2$ (STO-3G, CAS(12,12)) Hamiltonians via
  PySCF; apply \texttt{measurement\_grouping\_v3.py} unchanged; report
  $M_{\mathrm{opt}}$ ungrp / QWC / FC + reductions. Flag any case where
  FC ceases to dominate.

\item \textbf{Noise robustness of the FC$>$QWC advantage.}
  \emph{Basis.} FC groups need a joint Clifford diagonalization circuit,
  adding gate depth; on noisy hardware the added error could offset the
  sample-cost win. Our replication was noise-free statevector; the review
  acknowledges but does not quantify this trade-off.
  \emph{Next steps.} On H$_2$, run QWC and FC measurement through a
  Qiskit fake-backend noise model (FakeManila, FakeGuadalupe) plus a
  parametrized readout-error channel (0.1--5\%). Measure realized MSE per
  total shot vs.\ noise level; find the crossover noise strength at which
  QWC beats FC despite a worse noiseless bound.

\item \textbf{ML-driven shot-budget allocation vs.\ the $\sqrt{\text{Var}}$
  optimum under noisy variance estimates.}
  \emph{Basis.} $M_{\mathrm{opt}}$ assumes exact $\mathrm{Var}(\hat H_\alpha)$.
  In practice variances are estimated from a small pilot batch and are
  themselves noisy; naive plug-in of $\sqrt{\widehat{\text{Var}}_\alpha}$ can
  waste shots. This mismatch is treated in the shadow literature but not
  deeply in the review's grouping sections.
  \emph{Next steps.} Two-stage protocol on LiH FC groups (35 arms): pilot
  batch of $K \in \{10, 50, 200\}$ shots per fragment; then allocate
  remaining budget by plug-in $\sqrt{\widehat{\text{Var}}}$. Compare realized
  MSE to (i) known-variance $\sqrt{\text{Var}}$ allocation and (ii) a
  Thompson-sampling bandit over per-fragment variance posteriors. Report
  the pilot-$K$ at which the bandit stops helping.
\end{enumerate}
